\documentclass[11pt]{article}
\usepackage[margin=1in]{geometry}
\usepackage{amsmath,amssymb}
\usepackage{enumitem}
\usepackage{bookmark}
\hypersetup{
  pdftitle={PSET 8: Multivariate distributions},
  hidelinks}

\title{PSET 8: Multivariate distributions}
\author{}
\date{\vspace{-2.5em}}

% Problems are numbered 1-5 across the document.
\newcounter{problem}
\newcommand{\problem}[1]{\refstepcounter{problem}\section*{Problem \theproblem: #1}}

\begin{document}
\maketitle

% Shared assignment questions (header block).
% Include at the top of any problem set with:  % Shared assignment questions (header block).
% Include at the top of any problem set with:  % Shared assignment questions (header block).
% Include at the top of any problem set with:  \input{assignment-qs}
\section*{Assignment Qs}

\begin{itemize}
\item Name
\item How long did this problem set take you (round to nearest hour)? \quad
  0-1 hour \ 2-3 hours \ 4-5 hours \ 5+ hours
\item How difficult was this problem set? \quad
  very easy \ 1 \ 2 \ 3 \ 4 \ 5 \ very challenging
\end{itemize}

\section*{Assignment Qs}

\begin{itemize}
\item Name
\item How long did this problem set take you (round to nearest hour)? \quad
  0-1 hour \ 2-3 hours \ 4-5 hours \ 5+ hours
\item How difficult was this problem set? \quad
  very easy \ 1 \ 2 \ 3 \ 4 \ 5 \ very challenging
\end{itemize}

\section*{Assignment Qs}

\begin{itemize}
\item Name
\item How long did this problem set take you (round to nearest hour)? \quad
  0-1 hour \ 2-3 hours \ 4-5 hours \ 5+ hours
\item How difficult was this problem set? \quad
  very easy \ 1 \ 2 \ 3 \ 4 \ 5 \ very challenging
\end{itemize}


\problem{Properties of a joint PDF}\label{prob:jointpdf}

Continuous random variables \(X\) and \(Y\) have the following joint
probability density function (PDF):\footnote{Inspired by Grimmer HW12.1}
\[
\begin{aligned}
f_{XY} (x, y) &= \begin{cases}
k x^3 y^2 & \mbox{where} \ 0 < x, y < 6 \\
0 & \mbox{otherwise}
\end{cases}
\end{aligned}
\]

Note: \(0 < x,y < 6\) means that both \(x\) and \(y\) are between 0 and
6; it does not mean that \(x\) is greater than 0 and \(y\) is less than
6.

\begin{enumerate}[label=\alph*.]
\item Find \(k\).
\item Find the marginal PDF of \(X\), \(f_X (x)\).
\item Find the marginal PDF of \(Y\), \(f_Y (y)\).
\item Find \(\Pr(X < 2, \ Y > 4)\).
\end{enumerate}

\problem{Calculating conditional PDF}\label{prob:condpdf}

Let \(f(x,y) = 18x^2y^2\) for \(0 \leq x \leq y \leq
1\).\footnote{Inspired by Grimmer HW12.4}

\begin{enumerate}[label=\alph*.]
\item Sketch the region on which \(f(x,y)\) is positive.
\item Find \(f(x|y)\).
\item Are \(X\) and \(Y\) independent? Explain your reasoning, paying
  attention to the region on which \(f(x,y)\) is defined.
\end{enumerate}

\problem{Moments and independence}

Continuous random variables \(W\) and \(Z\) have the following joint
probability density function (PDF):\footnote{Inspired by Grimmer HW12.1}
\[
\begin{aligned}
f_{WZ} (w, z) &= \begin{cases}
k w^2 z^3 & \mbox{where} \ 0 < w, z < 2 \\
0 & \mbox{otherwise}
\end{cases}
\end{aligned}
\]

\begin{enumerate}[label=\alph*.]
\item Find \(k\).
\item Find E\([W]\).
\item Find \(Var(W)\).
\item Find \(Cov(W, Z)\).
\item Are \(W\) and \(Z\) independent? Explain your reasoning using
  mathematical concepts from the course.
\item What is the PDF of \(W\) conditional on \(Z\), \(f_{W|Z} (w|z)\)?
\end{enumerate}

\problem{Properties of joint random variables}

Suppose the following:\footnote{Inspired by Grimmer HW12.3}

\begin{itemize}[nosep]
\item \(E[D] = 8\)
\item \(E[F] = 4\)
\item \(E[DF] = 10\)
\item \(Var(D) = 30\)
\item \(Var(F) = 60\)
\end{itemize}

\begin{enumerate}[label=\alph*.]
\item What is \(Cov(D,F)\)?
\item What is the correlation between \(D\) and \(F\)?
\item What is \(Var(D + F)\)?
\item What is \(Var(D - F)\)?
\item Suppose you multiplied \(F\) by 2 to generate a new variable,
  \(H\). What is \(Cov(D,H)\)?
\item What is \(Cor(D,H)\)? How does this compare to your answer to
  part (b) of this question?
\item Suppose instead that \(Var(D) = 40\). How would this change
  \(Cor(D,F)\)?
\end{enumerate}

\problem{Continuous Bayes' theorem}

Previously, we used Bayes' theorem to link the conditional probability
of discrete events \(A\) given \(B\) to the probability of \(B\) given
\(A\). There is an analogous Bayes' theorem that relates the
conditional densities of random variables \(X\) and
\(\theta\):\footnote{Inspired by Grimmer HW12.5}
\[
f(\theta \mid X) = \dfrac{f(X \mid \theta) f(\theta)}{\int f(X \mid \theta) f(\theta) \, d\theta}
\]

\begin{enumerate}[label=\alph*.]
\item Return to the joint PDF from Problem~\ref{prob:condpdf}. Using
  the conditional density \(f(x|y)\) and the marginal \(f_Y(y)\) you
  worked with there, apply Bayes' theorem to find \(f(y|x)\). Here
  \(y\) plays the role of \(\theta\) and \(x\) plays the role of the
  observed data \(X\). Be careful with the limits of integration in
  the denominator. 
\item Now compute \(f(y|x)\) directly as \(f(x,y) / f_X(x)\) and
  confirm that it matches your answer to part (a). What familiar
  quantity does the denominator of Bayes' theorem turn out to be here?
\end{enumerate}

% Shared AI and Resources statement.
% Include in any problem set with:  % Shared AI and Resources statement.
% Include in any problem set with:  % Shared AI and Resources statement.
% Include in any problem set with:  \input{ai-statement}
\section*{AI and Resources statement}

Please list (in detail) all resources you used for this assignment. If
you worked with people, list them here as well. It is not enough to say
that you used a resource for help, you need to be specific on the link
and \emph{how} it was helpful. W/R/T gen AI tools (including GPT,
Claude, etc.) you cannot use them to do work on your behalf---you
cannot put in any of the questions, etc. You can ask for help on logic
/ sample problems. If you do use GPT or other AI tools, you need to
provide a link to your chat transcript. Any suspected academic
integrity violations will be immediately reported to the Dean of
Students.

\section*{AI and Resources statement}

Please list (in detail) all resources you used for this assignment. If
you worked with people, list them here as well. It is not enough to say
that you used a resource for help, you need to be specific on the link
and \emph{how} it was helpful. W/R/T gen AI tools (including GPT,
Claude, etc.) you cannot use them to do work on your behalf---you
cannot put in any of the questions, etc. You can ask for help on logic
/ sample problems. If you do use GPT or other AI tools, you need to
provide a link to your chat transcript. Any suspected academic
integrity violations will be immediately reported to the Dean of
Students.

\section*{AI and Resources statement}

Please list (in detail) all resources you used for this assignment. If
you worked with people, list them here as well. It is not enough to say
that you used a resource for help, you need to be specific on the link
and \emph{how} it was helpful. W/R/T gen AI tools (including GPT,
Claude, etc.) you cannot use them to do work on your behalf---you
cannot put in any of the questions, etc. You can ask for help on logic
/ sample problems. If you do use GPT or other AI tools, you need to
provide a link to your chat transcript. Any suspected academic
integrity violations will be immediately reported to the Dean of
Students.


\section*{(Optional --- complete elsewhere) Submission of practice questions}

Submit practice questions for the final exam here:
\url{https://forms.gle/sskPerZjthPYPbVt7} Note that we need at least 10
people to submit before there's enough to circulate!

\end{document}